\documentclass[11pt]{amsart}

\usepackage[T1]{fontenc}
\usepackage{lmodern}
\usepackage{microtype}
\usepackage{amsmath,amssymb}
\usepackage[hidelinks]{hyperref}

\newtheorem{theorem}{Theorem}
\newtheorem{lemma}[theorem]{Lemma}

\title[An elementary proof of Erlbacher's counterexample]
{An Elementary Proof of Erlbacher's Counterexample to TxGraffiti Conjecture~9}
\date{}

\hypersetup{
  pdftitle={An Elementary Proof of Erlbacher's Counterexample to TxGraffiti Conjecture 9}
}

\subjclass[2020]{05C50, 05C69}
\keywords{zero forcing number, domination number, cubic graph, fort, counterexample}

\begin{document}

\begin{abstract}
Erlbacher identified and computationally verified a connected cubic
triangle-free graph on $14$ vertices that refutes Conjecture~9 of Davila,
generated by TxGraffiti.  We give a self-contained elementary proof that its
domination number is $4$ and its zero forcing number is $7$.  Hence
$Z(G)=\gamma(G)+3$, contradicting the proposed bound
$Z(G)\leq\gamma(G)+2$ for connected cubic diamond-free graphs.  The lower
bound on $Z(G)$ follows from an explicit family of forts.
\end{abstract}

\maketitle

\section{The counterexample}

All graphs are finite, simple, and undirected.  A set $D\subseteq V(G)$ is
\emph{dominating} if $N[D]=V(G)$, and the minimum size of such a set is the
\emph{domination number} $\gamma(G)$.  Starting from a set of blue vertices,
the zero forcing rule allows a blue vertex with exactly one white neighbor to
force that neighbor blue.  The minimum size of an initial set that can color
all vertices blue is the \emph{zero forcing number} $Z(G)$; see
\cite{AIM2008}.

Davila posed the following as Conjecture~9 of \cite{Davila2024}, generated by
TxGraffiti:
\begin{equation}\label{eq:conjecture}
  Z(G)\leq \gamma(G)+2
\end{equation}
for every connected cubic diamond-free graph $G$.  Erlbacher
\cite{Erlbacher2026} exhibited the graph below and verified its parameters by
exhaustive computation.  The same artifact gives larger computational examples
in a chain family.  The contribution here is a noncomputational proof for the
$14$-vertex member.

Let
\[
\begin{aligned}
 A&=\{a_1,a_2,a_3\}, & B&=\{b_1,b_2,b_3\},\\
 C&=\{c_1,c_2,c_3\}, & D&=\{d_1,d_2,d_3\}.
\end{aligned}
\]
Define $G$ on $A\cup B\cup C\cup D\cup\{x,y\}$ by
\[
\begin{aligned}
E(G)={}&\{a_i b_j:1\leq i,j\leq 3,\ (i,j)\neq(3,1)\}\\
&\cup\{c_i d_j:1\leq i,j\leq 3,\ (i,j)\neq(3,1)\}\\
&\cup\{a_3x,xb_1,c_3y,yd_1,xy\}.
\end{aligned}
\]
Equivalently, take two copies of $K_{3,3}$, subdivide one edge in each, and
join the two subdivision vertices.

\begin{theorem}\label{thm:main}
The graph $G$ is connected, cubic, and triangle-free, and
\[
  \gamma(G)=4
  \qquad\text{and}\qquad
  Z(G)=7.
\]
Consequently, $G$ is diamond-free and refutes \eqref{eq:conjecture}.
\end{theorem}

A nonempty set $F\subseteq V(G)$ is a \emph{fort} if every vertex outside
$F$ has either zero or at least two neighbors in $F$.

\begin{lemma}\label{lem:fort}
Every zero forcing set intersects every fort.
\end{lemma}

\begin{proof}
If a zero forcing set $S$ were disjoint from a fort $F$, then immediately
before the first force into $F$, the forcing vertex would lie outside $F$ and
have exactly one neighbor in $F$, contrary to the definition of a fort.
\end{proof}

\begin{proof}[Proof of Theorem~\ref{thm:main}]
The construction makes $G$ connected and cubic.  Each subdivided copy of
$K_{3,3}$ is triangle-free, and the edge $xy$ is a bridge, so $G$ is
triangle-free.  Since every diamond contains a triangle, $G$ is
diamond-free.

The set $\{a_3,b_1,c_3,d_1\}$ dominates $G$, so $\gamma(G)\leq4$.  Since
$G$ has $14$ vertices and maximum degree $3$, a set of $k$ vertices dominates
at most $4k$ vertices.  Hence $\gamma(G)\geq\lceil14/4\rceil=4$, and therefore
$\gamma(G)=4$.

For the zero forcing upper bound, let
\[
  S_0=\{y,a_1,a_2,b_2,c_1,c_2,d_2\}.
\]
The sequence
\[
 b_2\to a_3,\quad
 d_2\to c_3,\quad
 c_3\to d_3,\quad
 c_1\to d_1,\quad
 y\to x,\quad
 a_3\to b_3,\quad
 a_1\to b_1
\]
is valid, so $Z(G)\leq7$.

It remains to prove $Z(G)\geq7$.  Put $U=A\cup B$ and $W=C\cup D$.  The
following six subsets of $U$ are forts:
\[
 \{a_1,a_2\},\qquad \{b_2,b_3\},
\]
and
\[
 \{a_i,a_3,b_1,b_j\},
 \qquad i\in\{1,2\},\quad j\in\{2,3\}.
\]
For either two-element set, each outside vertex adjacent to it has exactly
two neighbors in it.  For a four-element set, the only outside vertices with
neighbors in the set are $x$, the other vertex of $\{a_1,a_2\}$, and the
other vertex of $\{b_2,b_3\}$; each has exactly two such neighbors.

No two vertices meet all six forts.  Indeed, any two-vertex set meeting the
first two forts must be $\{a_i,b_j\}$ with $i\in\{1,2\}$ and
$j\in\{2,3\}$, but it misses the four-element fort obtained from the other
choices of $i$ and $j$.  Lemma~\ref{lem:fort} therefore gives
\[
  |S\cap U|\geq3
\]
for every zero forcing set $S$.  By symmetry, $|S\cap W|\geq3$.

Suppose that $S$ is a zero forcing set of size six.  Then
\[
 |S\cap U|=|S\cap W|=3
 \qquad\text{and}\qquad
 S\cap\{x,y\}=\varnothing.
\]
Define
\[
\begin{aligned}
 \mathcal P&=\bigl\{\{a_1,a_3\},\{a_2,a_3\},
                    \{b_1,b_2\},\{b_1,b_3\}\bigr\},\\
 \mathcal Q&=\bigl\{\{c_1,c_3\},\{c_2,c_3\},
                    \{d_1,d_2\},\{d_1,d_3\}\bigr\}.
\end{aligned}
\]
For every $p\in\mathcal P$ and $q\in\mathcal Q$, the set
\[
  \{x,y\}\cup p\cup q
\]
is a fort.  If $p=\{a_i,a_3\}$, the other vertex of $\{a_1,a_2\}$ has
zero neighbors in $p\cup\{x\}$, while $b_1,b_2,b_3$ each have two.  If
$p=\{b_1,b_j\}$, the other vertex of $\{b_2,b_3\}$ has zero, while
$a_1,a_2,a_3$ each have two.  The verification for $q$ is symmetric.

Since $S$ misses $x$ and $y$, Lemma~\ref{lem:fort} implies that $S$ meets
$p$ or $q$ for every $p\in\mathcal P$ and $q\in\mathcal Q$.  Hence $S$ meets
every member of $\mathcal P$, or it meets every member of $\mathcal Q$.
In the first case, together with the two two-element forts above,
$S\cap U$ meets all six pairs
\[
 \{a_1,a_2\},\ \{a_1,a_3\},\ \{a_2,a_3\},\qquad
 \{b_1,b_2\},\ \{b_1,b_3\},\ \{b_2,b_3\}.
\]
These are the edges of two vertex-disjoint copies of $K_3$, so at least two
vertices from each copy are required.  Thus $|S\cap U|\geq4$, a
contradiction.  The second case similarly gives $|S\cap W|\geq4$.
Therefore no six-vertex set is zero forcing, so $Z(G)\geq7$.  Together with
the forcing sequence, this proves $Z(G)=7$.
\end{proof}

The edge $xy$ is a bridge, but bridgelessness is not a hypothesis of
Conjecture~9.  Moreover, $x$ and $y$ are centers of induced claws, so the
example does not affect Davila's theorem for connected cubic claw-free
graphs.

\begin{thebibliography}{9}

\bibitem{AIM2008}
AIM Minimum Rank--Special Graphs Work Group,
\emph{Zero forcing sets and the minimum rank of graphs},
Linear Algebra Appl. \textbf{428} (2008), no.~7, 1628--1648.
\href{https://doi.org/10.1016/j.laa.2007.10.009}
{doi:10.1016/j.laa.2007.10.009}.

\bibitem{Davila2024}
R.~R. Davila,
\emph{Another conjecture of TxGraffiti concerning zero forcing and domination
in graphs},
arXiv:2406.19231v2, 2024.
\href{https://doi.org/10.48550/arXiv.2406.19231}
{doi:10.48550/arXiv.2406.19231}.

\bibitem{Erlbacher2026}
J.~Erlbacher,
\emph{Demonstrandum---wave-1 verified artifacts},
version~3, Zenodo, July~8, 2026.
\href{https://doi.org/10.5281/zenodo.21269439}
{doi:10.5281/zenodo.21269439}.
See \texttt{problems/p2-factory/kills/davila-conj9}.

\end{thebibliography}

\end{document}